\section{Selected shell observations and geometric limits}\label{sec:geometry221}
This section groups one positive local construction with two sharp limits
on extending it. All energies below use counting measure and original
positive-axis starting-site coefficients. Write
\[
 V_Y(f)=\sum_Y|f-\overline f_Y|^2,\qquad
 E_{a,Y}(f)=\sum_{x,x+e_j\in Y}a(x)|f(x+e_j)-f(x)|^2.
\]
No coefficient reflection or endpoint averaging is implicit. Normalized
counting means are $q_Y(f)=|Y|^{-1/2}\sum_Y f$. Dividing every term by
the ambient volume gives the probability-normalized convention used earlier.
The local construction does not cover the whole domain. The first limit
concerns a smaller internal-cell certificate; the second concerns an actual
annular graph. Neither is a counterexample to the unrestricted full-domain
question in Section~27.

\subsection{Coefficient-selected unit shells with disjoint faces}
\begin{theorem}\label{thm:selected221}
Let $D\ge3$, $d=D-1$, $P=D/2$, $p=d/2$, $q=2d/(d+2)$,
$L\ge1$ an integer, $z\in\mathbb Z^D$, and $\mu>0$. In the nonwrapping integer lattice, on
$\Omega=\{x\in\mathbb Z^D:L\le\|x-z\|_\infty\le2L-1\}$ let $a(x)>0$ and put
$B=\sum_\Omega(\mu/a)^P$. There is a coefficient-only selected
$k\in\{L,\ldots,2L-1\}$ whose shell $A_k=\{x\in\mathbb Z^D:\|x-z\|_\infty=k\}$ has exactly $2D$
disjoint rectangular faces $F$, such that for every real $f$
\begin{equation}\label{eq:selected221}
 \sum_F V_F(f)\le \frac{C_DB^{2/D}}{\mu}\sum_F E_{a,F}(f)
 \le \frac{C_DB^{2/D}}{\mu}E_{a,A_k}(f).
\end{equation}
Here $K_2=2$ and, for $d\ge3$,
\[
 K_d=\frac{2\alpha_d d^{1/q'}}{d(1-2^{-1/d})},\quad
 \alpha_d=\frac{2(d-1)}d,\quad q'=\frac{2d}{d-2},\qquad
 C_D=4K_d^2 3^{2d/q}d^{2/d}.
\]
In particular $C_3=2592$. If $cC_DB^{2/D}/\mu\le1$, these actual
means give
\[
 E_{a,A_k}(f)+c\sum_F|q_F(f)|^2\ge c\sum_{A_k}|f|^2.
\]
The same assertion holds coordinatewise without an additional vector-dimension
factor. It supplies local rank $2D$ per scalar coordinate only.
\end{theorem}
\begin{proof}
Set $b=\mu/a$, $W_j=\sum_{A_j}b^P$ and choose the least minimizer of
$W_j$. Disjoint radial layers give $W_k\le B/L$. Relative to $z$, take
\[
 F_{i,\sigma}=[-k+1,k-1]^{i-1}\times\{\sigma k\}
                 \times[-k,k]^{D-i},\qquad \sigma\in\{-1,1\}.
\]
The first coordinate attaining absolute value $k$ assigns each shell site
to a unique face. Each tangential side has $2k-1$ or $2k+1$ sites;
the faces are nonempty, including $k=1$, and $|F|\le(4L)^d$.

We spell out the adapter to the complete counting cube Sobolev proof
retained earlier in this manuscript. Fold the $(2k+1)$-site interval
onto each $(2k\pm1)$-site tangential interval by the endpoint-repeating
walk $0,1,\ldots,\ell-1,\ell-1,\ldots,0,0,\ldots$.
Every site has between one and three preimages. Every internal edge has
at most three preimages: for $k\ge2$ the walk takes at most two steps
beyond its first traversal, while for $k=1$ the possible lengths are
one and three. The constant one-site map has no nonconstant edges.
The product fold therefore multiplies an unweighted edge sum by at most
$3^d$. For its extension $\widetilde f$, minimizing over constants gives
$V_F(f)\le V_{Q_{2k+1}}(\widetilde f)$. Thus the proved cube inequality
with exponent $q$ gives
\[
 V_F(f)\le K_d^2 3^{2d/q}
                 \left(\sum_{e\subset F}|df(e)|^q\right)^{2/q}.
\]
This fold is entirely unweighted; reversing an edge transports no
coefficient. Apply H\"older next on the original face edges:
\[
 \left(\sum_{e\subset F}|df|^q\right)^{2/q}
 \le\frac{d^{1/p}}\mu E_{a,F}(f)
               \left(\sum_F b^p\right)^{1/p}.
\]
Every site starts at most $d$ such edges. Since $p<P$,
\[
 \left(\sum_Fb^p\right)^{1/p}
 \le |F|^{1/p-1/P}W_k^{1/P}
 \le(4L)^{2/D}(B/L)^{2/D}\le4B^{2/D}.
\]
Here $d(1/p-1/P)=2/D$ is the radial cancellation. Summing over faces
uses disjoint retained edges. Finally
$\sum_{A_k}|f|^2=\sum_FV_F(f)+\sum_F|q_F(f)|^2$ proves the penalty
assertion. All starting sites, including incoming ones, occur in the
actual edge sums; no local moment has been assumed.
\end{proof}

The scale cannot be improved to $o(L^2)/\mu$ on homogeneous profiles.
Indeed $a=\mu$ gives $B=(4L-1)^D-(2L-1)^D\le(4L)^D$.
For every candidate radius and $f(x)=x_2-z_2$, the face $F_{1,+}$
alone has variance $(2k+1)^d((2k+1)^2-1)/12$, whereas
$E_{a,A_k}(f)\le\mu|A_k|\le2D\mu(2k+1)^d$. Consequently
\[
 \frac{\sum_FV_F(f)}{E_{a,A_k}(f)}
 \ge\frac{L^2}{6D\mu}\ge\frac{B^{2/D}}{96D\mu}.
\]
This uses the whole shell energy, including its face interfaces. It is
an all-candidate-radii lower scale bound. A unit shell has only order
$L^{D-1}$ sites: its $2D$ observations do not cover the unselected
layers or interior. Repeating independently on order $L$ layers costs
order $L$ rows. This independent-layer repetition alone does not supply
a global covering and charging theorem at the original target count.
Section~\ref{sec:balanced223} derives the finite scalar construction by
a different balanced-ring mechanism.

\subsection{Sharp rank cost of discarding canonical interfaces}
\begin{theorem}\label{thm:filament221}
Fix $D\ge3$, integer $s\ge1$, $R=s^D$, $n=2R+1$, and
$X=\{-R,\ldots,R\}^D$, $N=n^D$. Let $a=A=s^2$ on the $2D$
punctured coordinate half-axes and $a=H=n^2$ elsewhere. Its actual
mean $\mu$ satisfies $1\le a\le2\mu$, $\mu\ge1$, and
\[
 1\le\frac1N\sum_X(\mu/a)^{D/2}<2D+1,\qquad
 \max\{1,(n/\sqrt{1+\mu})^D\}\le2^{D/2}.
\]
For the actual full energy one rectangular mean suffices:
\begin{equation}\label{eq:fullfil221}
 E_{a,X}(f)+c_D|q_X(f)|^2\ge c_D\|f\|_2^2,
 \qquad c_D=[64D^2(2D+1)]^{-1}.
\end{equation}
For any coordinate permutation $\pi$, partition instead by the first
nonzero coordinate, with the origin as a singleton. Let $P_0$ be its
$2D+1$ normalized means and $E_{\mathrm{int},\pi}$ its retained
internal-cell energy. Set $\ell=32D(D+1)s$ and
$m_s=\lfloor(R-1)/(4\ell)\rfloor$. If real-linear maps $C,V$ and
any $\Phi$ with $\Phi(0)=0$ obey, for all $f$ or near zero,
\[
 E_{\mathrm{int},\pi}(f)+K\|P_0f\|^2+\Phi(Cf)+\langle f,Vf\rangle
 \ge c_D\|f\|^2,
\]
then $\rank C+\rank V\ge2m_s$. No symmetry or positivity of $V$
is needed, and $K$ may be any finite scalar. This growing repair cost
has matching order $\Theta_D(s^{D-1})$ by an explicit refinement.
\end{theorem}
\begin{proof}
There are $w=2DR$ weak sites. With $P=D/2$,
\[
 \mu=H-\frac{w(H-A)}N,\qquad S:=\sum_Xa^{-P}=2D+1-w/N.
\]
Since $N\ge2R3^{D-1}$, $w/N\le D/3^{D-1}\le1/3$ and
$\mu\ge2H/3$. The stated moment equals $(\mu/H)^PS<2D+1$;
its lower bound is Jensen. The target factor follows from $\mu\ge H/2$.
The earlier complete cube proof and original-start H\"older give
\[
 V_X(f)\le K_D^2(DS)^{2/D}E_{a,X}(f)
          \le64D^2(2D+1)E_{a,X}(f).
\]
For clarity, its constant satisfies $K_D\le8\sqrt D$:
$D(1-2^{-1/D})\ge1/2$ follows from
$(1-1/(2D))^D\ge1/2$, while $\alpha_D\le2$ and
$D^{1/q'}\le\sqrt D$. The variance identity proves
\eqref{eq:fullfil221}. Periodic wrap edges, if added, only help it.

The canonical cells are
\[
 B_{i,+}=\{x:x_{\pi(j)}=0\ (j<i),\ 1\le x_{\pi(i)}\le R\},
\]
with all later coordinates free in $[-R,R]$; define $B_{i,-}$ using
$[-R,-1]$ and include $\{0\}$. These disjoint integer rectangles
cover $X$. Their final pair are weak half-axes for every $\pi$.
For $R\ge2$, the mean-zero linear field $f_t=t-(R+1)/2$ on one final half-axis,
zero elsewhere, has
\[
 P_0f=0,\qquad
 \frac{E_{\mathrm{int},\pi}(f)}{\|f\|^2}
       =\frac{12s^2}{R(R+1)}\longrightarrow0.
\]
In the actual energy, each transverse incoming edge starts at a
two-nonzero-coordinate site and carries $H$, not $A$. Therefore any
field supported on the final two half-axes has
$E_{a,X}(f)\ge(D-1)H\|f\|^2$. The omitted interfaces explain the
certificate failure without contradicting the full-energy bound.

To obtain the rank obstruction, use the signed tent on $0\le t\le4\ell$:
\[
 g(t)=\begin{cases}t&0\le t\le\ell,\\
 2\ell-t&\ell\le t\le3\ell,\\t-4\ell&3\ell\le t\le4\ell.
 \end{cases}
\]
It has zero endpoints and sum, and
\[
 \sum g^2=2\ell(2\ell^2+1)/3,\qquad\sum|dg|^2=4\ell.
\]
Place $m_s$ copies on consecutive intervals
$[1+4\ell j,1+4\ell(j+1)]$ of each final half-axis, indexed by its
absolute coordinate. Nonzero supports and edge differences do not
overlap; shared endpoints are zero. Their span $G_s\subset\ker P_0$
has dimension $2m_s$, and for every $f\in G_s$
\[
 E_{\mathrm{int},\pi}(f)=\rho_s\|f\|^2,\qquad
 \rho_s=\frac{6s^2}{2\ell^2+1}
 \le\frac3{[32D(D+1)]^2}\le c_D/4.
\]
The last step is $4(D+1)^2\ge3(2D+1)$. If combined repair rank
is smaller than $2m_s$, rank-nullity supplies nonzero
$f\in G_s\cap\ker C\cap\ker V$. All repair terms vanish;
scaling toward zero also treats a neighborhood inequality. This
contradicts the prescribed lower bound. The statement is vacuous when
$m_s=0$, and asymptotically $m_s$ has order $s^{D-1}$.

For the upper repair order, split each canonical cell's strong complement
by its first nonzero later coordinate. There are $2(D-i)$ strong
rectangles per $B_{i,\pm}$, hence $2D(D-1)$ in total, each with
coefficient $H=n^2$ and largest side at most $n$. Split each of the
$2D$ weak axes into $R/s=s^{D-1}$ intervals of length $s$ and retain
the origin. The total is
\[
 J_s=1+2D(D-1)+2Ds^{D-1}.
\]
The proved interval-product inequality
$V_B\le L_B(L_B-1)\sum_{e\subset B}|df|^2/4$ gives
$E_{a,B}\ge4V_B$ for every nonconstant cell. Singletons have zero
variance. Summing and using the normalized mean identity gives
coercivity with penalty $c_D$ on these $J_s$ means even for the refined
internal energy. That energy is at most $E_{\mathrm{int},\pi}$, so
these rows also repair the original certificate (take
$\Phi(Cf)=c_D\|P_{\mathrm{ref}}f\|^2$ and a nonnegative base penalty).
This proves the matching order without an optimal multiplicative constant.
\end{proof}

This theorem stops transferring a shell gap to these canonical full-volume
cells while dropping their interfaces for free. It applies to every order
selected from $a$, but not to every rectangular partition. Indeed its
one-cell full-energy solution is explicit. The scalar tents have not been
proved coexact or divergence constrained.

\subsection{An actual off-center annular slowdown of sharp order}
\begin{theorem}\label{thm:offcenter221}
For $D\ge3$ and integer $s\ge2$, put $d=D-1$, $L=s^D$, $Q=4L+1$,
\[
 X=\{0,\ldots,4L\}^D,\quad I=\{1,\ldots,2L\}^D,
 \quad\Omega=X\setminus I,\quad N=Q^D-(2L)^D.
\]
On the plate $F=\{x_1=0,\ 1\le x_j\le2L\ (j>1)\}$ set
$a=A=s^{2d}$, and set $a=H=Q^2$ elsewhere in $\Omega$.
Its actual mean obeys $\mu\ge1$, $1\le a\le2\mu$, and
\[
 1\le M_a:=N^{-1}\sum_\Omega(\mu/a)^{D/2}
           <2(1+2^d).
\]
Nevertheless the best Poincar\'e constant for the actual nonwrapping
annular energy satisfies
\[
 \frac{s^2}{12d}\le
 \sup_{f\ \text{nonconstant}}\frac{V_\Omega(f)}{E_{a,\Omega}(f)}
 \le B_Ds^2,\qquad B_D=25J(1+8J),\quad J=2D.
\]
Thus the thickness-uniform centered critical-moment formula cannot extend
to arbitrary off-center holes without a quantitative geometric condition.
The full cube obtained by filling the hole admits one rectangular mean.
\end{theorem}
\begin{proof}
With $w=(2L)^d$ and $P=D/2$, exact counting gives
\[
 \mu=H-w(H-A)/N,\qquad
 S:=\sum_\Omega a^{-P}=2^d+(N-w)/Q^D.
\]
Indeed $A^P=L^d$. Since
$w/N\le[2(2^D-1)L]^{-1}\le1/14$, $\mu\ge13H/14\ge H/2$.
Jensen gives $M_a\ge1$, and
\[
 M_a=\mu^P S/N
 \le\frac{1+2^d}{1-2^{-D}}<2(1+2^d).
\]
Let $h(t)=\min(t-1,2L-t)$ on $1\le t\le2L$ and put
$f(0,x_2,\ldots,x_D)=\prod_{j>1}h(x_j)$ on $F$, zero elsewhere.
Direct finite sums give
\[
 \sum h=L(L-1),\qquad
 T:=\sum h^2=\frac{(L-1)L(2L-1)}3,\qquad
 \sum|dh|^2=2L-2.
\]
The normal incoming edge starts outside $X$, and the outgoing edge enters
$I$; both are absent. At every tangential plate boundary $h=0$, so all
strong incoming or outgoing boundary edges have zero difference.
No intercell edge in $\Omega$ is discarded. Consequently
\[
 \|f\|^2=T^d,\quad \sum f=[L(L-1)]^d,\quad
 E_{a,\Omega}(f)=Ad(2L-2)T^{d-1}.
\]
Cauchy on the $w$-site support gives $V_\Omega(f)\ge\|f\|^2/2$,
whereas $E_{a,\Omega}(f)/\|f\|^2=6dA/[L(2L-1)]\le6d/s^2$.
This proves the lower bound. Subtracting its mean yields a mean-zero
witness with the same energy and variance. Meanwhile the multiplier
in the centered formula is bounded independently of $s$:
\[
 Q^2(S/N)^{2/D}
 \le\left(\frac{1+2^d}{1-2^{-D}}\right)^{2/D}.
\]
Thus no dimension-only constant can supply that off-center formula.

For the matching upper bound, cover $\Omega$ by $J=2D$ actual slabs
\[
 S_{i,-}=\{x:x_i=0\},\qquad
 S_{i,+}=\{x:2L+1\le x_i\le4L\},
\]
with all other coordinates free in $[0,4L]$. Their volumes are
$v_-=Q^d$ and $v_+=2LQ^d$. Root a tree at $S_{1,+}$, attach every
slab with different normal to that root, and attach $S_{1,-}$ to
$S_{2,+}$. Paths have length at most two, join distinct normals,
and every parent is fat. Descendant volume is no larger than the
volume of any slab on its root path.

Write $E_i^0$ for a slab's unweighted internal energy, $m_i$ for its
mean, and $Z=\sum_iE_i^0$. The interval-product proof gives
$V_{S_i}\le Q^2E_i^0/4$. To control means without an inverse
overlap-width factor, for any nonempty subinterval $T_0$ of a
$Q$-site interval $B$, Jensen and path Cauchy give
\[
 |\overline g_B-\overline g_{T_0}|^2
 \le\frac1{Q|T_0|}\sum_{x\in B,y\in T_0}|g(x)-g(y)|^2
 \le Q\sum_{e\subset B}|dg(e)|^2.
\]
Averaging over identical fibers of $S_i$ in another slab's normal
direction gives
\[
 |m_i-\overline f_{S_i\cap S_j}|^2\le Q^2E_i^0/v_i.
\]
The restricted width may be one or $2L$. Applying this on both sides
of an intersection yields
$|m_i-m_j|^2\le2Q^2(E_i^0/v_i+E_j^0/v_j)$.
Along a root path use Cauchy, its length at most two and volume
monotonicity; each intermediate energy occurs at most twice. Hence
$v_i|m_i-m_{\rm root}|^2\le8Q^2Z$. The cover and variance identity
now give, for every real $g$,
\begin{align*}
 V_\Omega(g)&\le\sum_i\{V_{S_i}(g)+v_i|m_i-m_{\rm root}|^2\}\\
 &\le(1/4+8J)Q^2Z
 \le J(1/4+8J)Q^2 E_\Omega^0(g)
 \le B_Ds^2E_{a,\Omega}(g).
\end{align*}
Here every actual edge is counted at most $J$ times, $a\ge A$,
$Q\le5L$, and $L^2/A=s^2$. For $D=3$, $B_3=7350$.

Finally fill $I$ with coefficient $H$ and extend the same plate field
by zero. The exact identity is
\[
 E_{a,X}(f)=E_{a,\Omega}(f)+A\|f\|^2.
\]
The restored normal edge starts on the weak plate, so the coefficient
is $A$, not $H$. All other added differences are zero. On the full
cube $\sum_Xa^{-P}<1+2^d$, and the complete cube bound used above gives
\[
 V_X(g)\le64D^2(1+2^d)E_{a,X}(g).
\]
Thus $c_X=[64D^2(1+2^d)]^{-1}$ and the one mean $q_X$ supply
$E_{a,X}+c_X|q_X|^2\ge c_X\|g\|^2$ for every $g$.
\end{proof}

Exact centering has not been proved the weakest sufficient condition;
quantitative clearance might suffice. The annular quotient of order
$s^{-2}$ is a finite scalar Neumann quotient, not a physical mass gap.
The full-domain unrestricted question and coexact restrictions remain
untouched. An initial unreviewed candidate note incorrectly called the
restored edge strong and incoming. Its original and pre-freeze correction
are retained; the theorem and its regression fixture use the weak outgoing
start. No failed scientific execution was hidden or rerun to change that
criterion.

\subsection{Comparison, evidence and remaining question}
Radial selection, folds, finite Sobolev/H\"older, product tents and
rank-nullity are classical. The weighted method should be compared with
Andres--Deuschel--Slowik~\cite{ads221}, Section~3.1, Proposition~3.5 of arXiv:1306.2521v5: its conductance and cutoff setup is not this finite-face
construction. Birman--Solomyak's adaptive cube partitions~\cite{birman221} do not supply
an unsupported zero-smoothness partition theorem here.
Brudnyi~\cite{brudnyi221}, Sections~1--2 and Appendix~I, treats
function-dependent approximation on dyadic rings. Appendix~I, Remark~I.2,
also treats boundary-touching dyadic inner cubes; an exclusively strictly
interior description would be incomplete. Those classical results do not by
themselves verify the particular integer original-start weighted adapter
required here. Section~\ref{sec:balanced223} derives that adapter explicitly;
no claim that arbitrary shifted holes have uniform geometry is made. Exact priority of these quantitative specializations
is unestablished; this limited primary comparison is not exhaustive.

The frozen analytical source commits are
\begin{center}\scriptsize
\begin{tabular}{ll}
selected shell218 & \texttt{e040475e37bf6b7440625df41c723738dbe3dbd8}\\
internal certificate219 & \texttt{5298926f130f3d680dfaa633bbf82e780bd0af84}\\
off-center annulus220 & \texttt{0c79cd67a002c24a305b70cf707c420aa1bd9ae3}
\end{tabular}
\end{center}
Their complete source reports, objections, supplied context and original
finite transcripts are retained under the corresponding research folders.
All three received accept-with-limits source verdicts in a continuing,
shared-context Codex task; this is role-separated but not independent or
blinded review. Exact hidden backend identity is unattested. Source218
required only a citation locator correction after its freeze, preserved
as a bounded text delta; sources219 and220 required no mathematical
correction after review. These reports are distinct from this changed
PDF's required numerical assessment.

Actual owned Colab normal and optimized diagnostics agreed in each unit.
Source218 checked 160 fields, 960 face/fold cases, 192 moment cases,
20 scale cases and three rejected negative controls. Source219 checked
60 small-grid fields, 1680 refined-cell cases, six linear witnesses,
ten profiles, four low-space cases and three rejected negative controls.
Source220 checked 20 fields on four small geometries, 120 slab bounds,
200 directional transports, 120 root bounds, five fields on one full
weighted profile, five filled-cube cases, six closed-profile cases and
three rejected negative controls. These are selected finite rational
diagnostics, not all-size machine proofs, new Lean results or two-fresh-clone
terminal reproduction. Runtime allocation, real exits, source hashes and
release/readback records are preserved; billed compute units are unknown.

One positive local input has been derived from the coefficients, and two
specific extension routes have been stopped by proofs. These local results
alone did not supply a full-domain partition and count. Section~\ref{sec:balanced223}
now derives them for the original scalar-gradient question by a different
balanced-ring mechanism. Earlier stopped harmonic-cost, common-refinement,
canonical interface-discard and unconditional off-center routes remain stopped. PR29 remains at 20/41 and
TermSource $=0$. No interacting or continuum conclusion follows.
